% ============================================================
%  Chapter 22: GPU Fragment Shaders as Observation Apparatus
%  Source: oxford/publications/observation-framework/
%          mass-computing-observation-framework.tex
% ============================================================

\epigraph{%
  The shader evaluates the partition function.\\
  The partition function is the physical observable.\\
  Therefore the shader is a physical measurement.%
}{}

\section{From Triple Equivalence to GPU Observation}
\label{sec:ch22-gpu-obs}

The Triple Equivalence $\Osc \cong \Cat \cong \Part$ has a striking practical
consequence.
In the $\Part$ description, a physical measurement is the identification of a
partition cell.
A GPU fragment shader, when programmed to evaluate the partition depth field
$\Depth(\mathbf{x})$ on a pixel grid, is performing exactly this operation:
for each pixel (corresponding to a spatial point in the analyser field), it
identifies the partition cell and evaluates $\Depth$.
By $\Part \cong \Osc$, this is a physical observation of the ion's trajectory.

\begin{remark}
  This is not a simulation.
  A simulation approximates a physical process; the GPU shader \emph{is}
  the physical process, in the categorical sense.
  The distinction matters for the S-entropy of the result: a simulation
  accumulates simulation error; a partition observation accumulates only the
  Floor theorem's irreducible uncertainty $S_\mathrm{flat}(R) > 0$.
\end{remark}

\section{The Triple Observation Architecture}
\label{sec:ch22-triple}

The observation framework implements the triple equivalence as a dual-path
architecture:

\paragraph{Oscillatory path ($\Osc$).}
The Runge--Kutta integrator propagates the ion's classical trajectory through
the analyser field using the partition Lagrangian~\eqref{eq:partitionlagrangian}.
Output: time-series of position and momentum.

\paragraph{Partition path ($\Part$).}
The fragment shader evaluates the partition depth field at each point on a
spatial grid.
Output: spatial map of $\Depth(\mathbf{x})$.

\paragraph{Categorical path ($\Cat$).}
The S-entropy tracker monitors $(\Sk, \St, \Se)$ at each step, ensuring that
the observation respects the Floor theorem.
Output: S-entropy time-series.

The three paths converge to the same mass spectrum.
Discrepancy between the paths is an instrument calibration metric.

\section{Quality Metrics and Training Signal}
\label{sec:ch22-quality}

The dual-path interference --- the discrepancy between $\Osc$ and $\Part$
outputs --- provides a natural training signal for instrument calibration:
\begin{equation}
  \Delta_\mathrm{obs} = \|\Sk^\mathrm{osc} - \Sk^\mathrm{part}\|_2.
\end{equation}
A well-calibrated instrument has $\Delta_\mathrm{obs} \to 0$.
The Floor theorem bounds: $\Delta_\mathrm{obs} \geq S_\mathrm{flat}(R)$.

%%  SOURCE: integrate full derivations and panels from
%%          mass-computing-observation-framework.tex
%%  ADD: panel 1 (triple observation scheme), panel 2 (quality memory),
%%       panel 3 (training dual-path), panel 4 (throughput scaling).

\bigskip
\begin{tcolorbox}[colback=crownblue!5, colframe=crownblue!60!black,
                  title={\textbf{Chapter 22 Summary}}]
\begin{itemize}[noitemsep]
  \item By the Triple Equivalence, a GPU fragment shader evaluating the
    partition depth field is a physical observation apparatus.
  \item The triple observation architecture runs $\Osc$, $\Part$, and $\Cat$
    paths in parallel; all three converge to the same spectrum.
  \item Dual-path interference provides a natural calibration metric bounded
    below by the Floor theorem.
\end{itemize}
\end{tcolorbox}
